% !TeX root = ../main.tex
\section{Introduction}\label{sec:intro}

Every memory reference on this machine is served by one of four levels that differ by two orders of
magnitude in latency: a \qty{32}{\kibi\byte} L1d and \qty{256}{\kibi\byte} L2 per core, a
\qty{20}{\mebi\byte} shared L3, and DDR4 DRAM. Which level a kernel hits, and how many references it keeps in
flight, decides its performance far more than its instruction count. We measure this hierarchy on an Intel
i9-10900F with purpose-built microbenchmarks and hardware counters, interpret the results with AMAT and
Little's Law, and answer:
\begin{enumerate}[nosep]
  \item What is the zero-queue access latency of each level, and how do latency and bandwidth change with
        access pattern, access size and read/write mix? (\cref{sec:zero-queue,sec:pattern,sec:rw-mix})
  \item Where does added concurrency stop buying throughput and start costing latency, and does Little's
        Law hold with measured latency? (\cref{sec:intensity})
  \item How do working-set size, cache misses and TLB misses show up in counters and runtime?
        (\cref{sec:working-set,sec:cache-miss,sec:tlb})
  \item Which benchmark mistake produces misleading numbers? (\cref{sec:pitfall})
\end{enumerate}

Every number in the text is generated from the processed data (\code{data/summary.csv},
\cref{app:repro}); the conclusion (\cref{sec:conclusion}) summarises the findings.
